\documentclass[9pt,a4paper]{extarticle}

\usepackage[utf8]{inputenc}
\usepackage[T1]{fontenc}
\usepackage[spanish,es-nodecimaldot]{babel}
\usepackage{lmodern}
\usepackage[a4paper,margin=1.65cm]{geometry}
\usepackage{microtype}
\usepackage{enumitem}
\usepackage{xcolor}
\usepackage{hyperref}
\usepackage{titlesec}
\usepackage{fancyhdr}
\setlength{\parindent}{0pt}
\setlength{\parskip}{3pt}
\linespread{0.96}

\definecolor{azulugr}{HTML}{174A6E}
\hypersetup{
    colorlinks=true,
    linkcolor=azulugr,
    urlcolor=azulugr,
    pdftitle={Memoria de la Práctica 1: Carrito de la compra},
    pdfauthor={José Javier Torres Cantero}
}
\titleformat{\section}{\large\bfseries\color{azulugr}}{\thesection.}{0.55em}{}
\titlespacing*{\section}{0pt}{0.5em}{0.15em}
\setlist[itemize]{leftmargin=1.5em,itemsep=0pt,topsep=1pt,parsep=0pt}
\pagestyle{fancy}
\fancyhf{}
\fancyhead[L]{Desarrollo de Software}
\fancyhead[R]{Práctica 1}
\fancyfoot[C]{\thepage}
\setlength{\headheight}{13pt}

\title{\vspace{-2.0em}\LARGE\bfseries\color{azulugr} Carrito de la compra\\[0.2em]
       \large Aplicación web con Spring Boot y Thymeleaf}
\author{José Javier Torres Cantero\\[0.25em]\normalsize Práctica 1}
\date{27 de octubre de 2026}

\begin{document}
\maketitle
\vspace{-2.0em}
\begin{center}
\footnotesize\textit{Memoria descriptiva de arquitectura, funcionalidades y ejecución}
\end{center}

\section{Arquitectura}

La aplicación adopta una arquitectura web por capas. La clase
\texttt{CarritoApplication} inicia Spring Boot y activa el escaneo de componentes
del paquete principal. Los controladores MVC reciben las peticiones HTTP y
preparan el modelo para las vistas Thymeleaf; los servicios concentran las
operaciones y validaciones; y \texttt{ProductRepository}, basado en Spring Data
JPA, proporciona el acceso a la entidad persistente \texttt{Product}.

El modelo de producto almacena identificador, nombre y precio decimal. El
catálogo se conserva en una base de datos H2 en memoria. El servicio del carrito
guarda los identificadores de los productos en el ámbito de la sesión HTTP, por
lo que cada visitante mantiene su propio carrito mientras dura la sesión.

\section{Funcionalidades}

\begin{itemize}
    \item Página de inicio y catálogo público de productos.
    \item Alta, edición y eliminación de productos, restringidas al rol
    administrador.
    \item Carrito de compra por sesión, con operaciones para añadir y quitar
    artículos y volver al catálogo.
    \item Inicio y cierre de sesión mediante Spring Security, con control de
    acceso a las rutas de administración.
    \item Descarga del catálogo como sentencias SQL desde
    \texttt{/admin/export}.
\end{itemize}

\section{Tecnologías y configuración}

El proyecto utiliza Java 21, Spring Boot 4.1.2, Maven, Spring Web MVC,
Thymeleaf, Spring Data JPA, Hibernate, H2 y Spring Security. No depende de
Lombok ni de complementos adicionales del IDE: los constructores y métodos de
acceso están escritos en el código fuente.

La base de datos es volátil: sus datos se reinician al detener la aplicación.
El nombre inicial de administrador es \texttt{admin}; por seguridad, la
contraseña no está incluida en el repositorio. Antes de arrancar la aplicación
se debe definir \texttt{CARRITO\_ADMIN\_\allowbreak PASSWORD}. En Windows,
puede establecerse en la terminal con
\texttt{\$env:CARRITO\_ADMIN\_\allowbreak PASSWORD="una-clave-local"}; en Eclipse
o STS, se configura como variable de entorno de la ejecución. Opcionalmente,
\texttt{CARRITO\_ADMIN\_\allowbreak USERNAME} permite cambiar el usuario.

\section{Importación, ejecución y pruebas}

Se requiere un JDK 21. Para importar el proyecto en Eclipse o Spring Tool Suite
(STS), se descomprime el ZIP y se selecciona
\textit{File} $>$ \textit{Import} $>$ \textit{Maven} $>$ \textit{Existing Maven
Projects}; como directorio raíz se elige la carpeta \texttt{carrito},
que contiene \texttt{pom.xml}. Tras resolver las dependencias, en STS se puede
ejecutar \texttt{CarritoApplication.java} como \textit{Spring Boot App}; en
Eclipse sin Spring Tools, como \textit{Java Application}.

También se puede iniciar desde una terminal situada en el proyecto:
\begin{itemize}
    \item Windows: \texttt{.\textbackslash mvnw.cmd}
    \texttt{spring-boot:run}.
    \item Linux o macOS: \texttt{sh ./mvnw}
    \texttt{spring-boot:run}.
\end{itemize}
La aplicación queda disponible en
\url{http://localhost:8080}. Las pruebas automatizadas se ejecutan con
\texttt{.\textbackslash mvnw.cmd} \texttt{clean verify} en Windows o
\texttt{sh ./mvnw} \texttt{clean verify} en Linux y macOS. El proyecto se ha
validado con siete pruebas de persistencia, rutas, permisos, autenticación,
operaciones del carrito y exportación SQL.

La comprobación funcional incluye consultar \texttt{/products}, iniciar
sesión en \texttt{/login}, gestionar productos desde la administración,
verificar las operaciones del carrito en \texttt{/cart} y descargar la
exportación SQL. Las operaciones de escritura usan protección CSRF; al
verificarlas con Postman se debe conservar la sesión y enviar el token CSRF.

\end{document}
